\documentclass{article}

\usepackage[final]{neurips_2019}

\usepackage[utf8]{inputenc}
\usepackage[T1]{fontenc}
\usepackage{hyperref}
\usepackage{url}
\usepackage{booktabs}
\usepackage{amsfonts}
\usepackage{nicefrac}
\usepackage{microtype}
\usepackage{graphicx}
\usepackage{xcolor}
\usepackage{enumitem}
\setlist{leftmargin=*, topsep=2pt, itemsep=1pt, parsep=0pt}

\title{
  Same Teacher, Same Ideas: \\
  Distillation Homogenizes Content \\
  Across Distinct Base Models \\
  \vspace{1em}
  \small{\normalfont COMS 4705 Project Proposal}  \\
  \small{\normalfont \textbf{Keywords:} \textit{homogenization, synthetic data, knowledge distillation, content and style, open-ended generation}}
}

% One row for all three authors; the shared affiliation goes underneath.
\author{
  \begin{tabular}{c@{\hspace{3em}}c@{\hspace{3em}}c}
    Xuanyi Lyu & Jiajun Chen & Boshu Dou \\
    \normalfont\texttt{xl3693@columbia.edu} & \normalfont\texttt{jc6949@columbia.edu} & \normalfont\texttt{bd2847@columbia.edu} \\[2pt]
    \multicolumn{3}{c}{\normalfont Department of Computer Science, Columbia University}
  \end{tabular}
}

\begin{document}

\maketitle

\section{Key Information}

\begin{itemize}
    \item External collaborators: None.
    \item Mentor: We would like to request Tianyi Lorena Yan, Saif Punjwani, or Andrew Tang as our mentor.
    \item Sharing project: We are not sharing this project with another class or with research outside the course.
\end{itemize}

\section{Literature review and reflection}

\paragraph{Papers you read.}
Jiang et al.~\cite{jiang2025hivemind} introduce Infinity-Chat, a dataset of 26K real user queries that have many acceptable answers and no single correct one. They use it to show that one model gives similar answers across repeated samples and that different models give similar answers to each other, and they find the second effect to be the stronger one. They call this the Artificial Hivemind. Because the study compares released models, it describes the agreement between models without testing what causes it.

Fortier et al.~\cite{fortier2026convergence} ask where in training this homogeneity begins. In their controlled supervised fine-tuning (SFT) experiments, the fine-tuning data revealed and amplified convergence on targeted prompts but did not introduce it. Base models also showed a similar collapse when they were only prompted. The authors conclude that homogeneity is largely learned in pretraining and that SFT acts as a catalyst. Their SFT experiments vary a small instruction set on a single base model.

Cloud et al.~\cite{cloud2025subliminal} fine-tune a student on number sequences written by a teacher that has a trait such as a preference for owls. The student acquires the trait even when the data are filtered to remove any mention of it. The transfer disappears when the teacher and the student are built on different base models. Next to the first two papers, this result shows that teacher data carries more than its visible content and that what a student absorbs depends on its base model.

Two further papers shaped our design. Gudibande et al.~\cite{gudibande2023false} found that models fine-tuned on ChatGPT outputs copied its style far more than its factual ability, so a teacher effect may be mostly stylistic. Karouzos et al.~\cite{karouzos2026collapse} traced diversity loss through the post-training of Olmo~3 and found that the stage where diversity collapses depends on the composition of the training data.

\paragraph{Questions.}
Our initial question was whether the similarity between different models' open-ended answers comes from pretraining or from synthetic post-training data that the models share. Fortier et al. convinced us that this split cannot be settled at our scale, because we cannot change what a base model saw during pretraining. The teacher that writes the fine-tuning data is a factor we can set directly. Gudibande et al. then made style the main alternative explanation to rule out. Our final question is: when base models from different families are fine-tuned on synthetic data from the same teacher, do they converge in the content of their answers, and how much of that convergence remains once style is held fixed?

\paragraph{Positioning and motivating your idea.}
Evidence for homogeneity across models is observational~\cite{jiang2025hivemind,lee2025quantification}. Released models differ in pretraining data, post-training data and teachers at the same time, so the share due to the teacher cannot be separated. Controlled studies~\cite{fortier2026convergence,karouzos2026collapse,schaffelder2025eggs} vary the training data of one model or one model family and measure diversity inside that model. We keep the prompts and the training recipe fixed and change only the teacher. We then measure how similar students built on distinct base models become. Compared with Cloud et al., we study content that a reader can see, and we test transfer across base models, which is the setting where their hidden transfer did not occur. The question matters in practice because many open models are post-trained on outputs of a small number of strong models~\cite{lee2025quantification}. If students of one teacher converge in content, a shared teacher is a direct and controllable cause of the inter-model homogeneity reported by Jiang et al.

\section{Project description}
\label{sec:description}

\paragraph{Goal.}
We call a base model fine-tuned on one teacher's answers a student and test three hypotheses.
\begin{itemize}
    \item \textbf{H1.} Students built on different base models are closer to each other when they share a teacher, across three teachers and three base-model families.
    \item \textbf{H2.} Most of this effect remains after another model rewrites the answers into one neutral style.
    \item \textbf{H3.} When one teacher's training answers are rewritten in a second teacher's style, the resulting students match students of the first teacher in content and students of the second teacher in style.
\end{itemize}
H3 is the causal test of the content explanation, because it changes the style of the training data and leaves its content in place. A secondary goal is to measure how the effect grows with data size. Stretch goals are to test whether the convergence survives a second SFT stage on human-written answers, and whether released models distilled from one teacher onto different bases~\cite{deepseek2025r1} behave alike.

\paragraph{Pilot results.}
Brackets give 95\% bootstrap confidence intervals. We first sampled four hosted teachers on 100 Infinity-Chat prompts: DeepSeek-V4.1-Flash, Gemma-4-26B-A4B, GLM-5.2 and gpt-oss-120b. Their answers differed in content, with a mean distance of 0.62 [0.45, 0.82] on a calibrated scale where 0 is the distance created by paraphrasing one teacher's answers and 1 is the distance created by telling that teacher to avoid the typical answer. We then fine-tuned Qwen3-4B-Base~\cite{yang2025qwen3} and Olmo-3-7B~\cite{olmo2025olmo3} on about 4,900 answers from DeepSeek or Gemma with two seeds per cell, and sampled 30 answers per student on 100 further prompts. Sharing a teacher brought students of different bases closer by 71\% [66, 76] of the distance that switching teachers adds within one base. After gpt-oss-120b rewrote 11,536 student answers into plain prose, the share of formatted answers fell from 28--48\% to 10--12\%, and 79\% [68, 88] of the shared effect remained.

\paragraph{Task.}
We compare the answer sets that students give to open-ended prompts unseen in training. The two teachers read the held-out prompt ``What if the entirety of Season 8 was all in Gendry's head while rowing?'' differently. Gemma treats the season's events as symbols of Gendry's inner conflict, such as guilt or a search for identity. DeepSeek retells the events as things Gendry imagined and ends with him still rowing. We had GLM-5.2 label each answer by the reading it takes. The symbolic reading appears in 28 of Gemma's 30 answers and in none of DeepSeek's. Gemma's students on Qwen and Olmo take it in 90\% and 100\% of their answers, and DeepSeek's students in 38\% and 40\%.

\paragraph{Data.}
Training prompts are the 4,938 prompts from the open-ended categories of databricks-dolly-15k~\cite{conover2023dolly} that have no context passage and that every teacher answered completely. Teachers give one direct answer of at most 100 words, sampled independently for each student run. Evaluation prompts are drawn at random from the 26K Infinity-Chat pool after removing the 100 prompts that Jiang et al. selected for high convergence, and a keyword filter removes unsafe and context-dependent prompts. Final results use 200 new prompts from the untouched pool, apart from the 300 pilot prompts.

\paragraph{Methods.}
\textit{Students.} We fine-tune LoRA adapters~\cite{hu2022lora} with rank 16 and $\alpha=32$ for 600 steps at batch size 16 and learning rate $2\times10^{-4}$, using our own training loop. The main grid has 27 students: three base models from different families (Qwen3-4B-Base, Olmo-3-7B and SmolLM3-3B-Base~\cite{smollm3}), three teachers (DeepSeek, Gemma and GLM) and three seeds. The three teacher pairs span calibrated distances of 0.20 to 0.84 in the pilot, so the grid also tests whether the shared effect grows with teacher distance.
\textit{Style normalization.} Our primary rewriter gpt-oss-120b turns the first 15 answers per student and prompt into plain prose for every student, with instructions to keep every claim and example. As a robustness check, Llama 4 Maverick from a fifth model family repeats this on a subset stratified by base, teacher and prompt category.
\textit{Cross-style data.} For H3, Gemma rewrites each of DeepSeek's training answers in its own words and style with the content unchanged, and DeepSeek does the same for Gemma's answers. As a control, the same source answers are also rewritten by their own teacher. Cross-style and control students use Qwen and Olmo with two seeds each and train on every prompt whose rewrites are complete answers in every condition. We hand-check a random sample of these training rewrites, mixed with decoys, for content fidelity.
\textit{Data scaling.} At a fixed budget of 600 update steps, we train students on nested random sets of 500, 1,000 or 2,500 distinct examples for two bases and two teachers.

\paragraph{Evaluation.}
For two students and one prompt, we embed every answer with all-mpnet-base-v2~\cite{reimers2019sbert} and compute the energy distance~\cite{szekely2013energy} between their two answer sets. Energy distance compares whole answer distributions, so it also detects shifts in how often each idea is picked. Sets have 30 answers, and all metrics are averaged over prompts. Comparisons before and after normalization use the same first 15 answers per student and prompt, cut to one word budget per prompt.
\begin{itemize}
    \item \textbf{Shared-teacher effect} $S$: the mean distance between students on different bases with different teachers, minus the same quantity for students with the same teacher.
    \item \textbf{Carryover} $S/(T-N)$, where $T$ is the distance between students of different teachers on the same base and $N$ is the distance between seeds. A value of 1 means that sharing a teacher brings students of different bases as close as switching teachers separates students of one base.
    \item \textbf{Normalized retention} $R=S_{\mathrm{norm}}/S_{\mathrm{orig}}$, computed on the same answers. A normalized answer is source-matched if it is closer to its source than to the student's other answers to that prompt. The pilot's rate was 91\%, and the full study adds a stratified hand check of content fidelity.
    \item \textbf{Cross-base teacher identification}: a logistic-regression classifier on answer embeddings is trained to name the teacher on two bases and half of the prompts, then tested on the third base and the disjoint held-out prompts, before and after style normalization (chance $1/3$).
    \item \textbf{H3 index}: a cross-style student's distance to its style teacher's students minus its distance to its content teacher's students, scaled by the distance between the two groups. The reference students share the student's base but neither its training draw nor its LoRA seed. H3 predicts a negative index on surface style features such as length and formatting rate, and a positive index on embeddings after style normalization. Self-rewrite controls show the effect of rewriting alone.
\end{itemize}
We repeat the distance metrics with bge-small and TF-IDF. For human evaluation, all three of us label the same 150 blinded cross-base answer pairs as sharing the main idea or not. We report Fleiss' $\kappa$ and compare the same-idea rate between same-teacher and different-teacher pairs. Qualitatively, we will describe what students inherit on the prompts with the largest shared effect, as in the Task example.

\paragraph{Baselines.}
\begin{itemize}
    \item \textbf{Different-teacher students} on different bases are the counterfactual that defines $S$.
    \item \textbf{Human-data students} trained on Dolly's human responses to the same prompts separate the effect of one teacher from that of sharing any fixed dataset.
    \item \textbf{Untuned base models} prompted with three human-written examples in the training format measure the similarity that the bases bring from pretraining, which Fortier et al. see as the main source of homogeneity.
    \item \textbf{Seed noise} $N$ gives a floor, and the distance between two teachers gives a reference.
\end{itemize}
We compare every baseline with the main grid as written and after style normalization. We implement all baselines ourselves because no prior work reports a cross-base teacher effect.

\paragraph{Justification.}
\begin{itemize}
    \item \textbf{Relevance:} the project studies how fine-tuning data shapes the text that language models generate.
    \item \textbf{Technical component:} we implement a LoRA training loop and a cross-base teacher classifier.
    \item \textbf{Language data:} the models learn from and generate English text in response to real user prompts.
    \item \textbf{Feasibility:} the pilot runs end to end, and the study adds about 50 LoRA runs on 3B--7B models.
    \item \textbf{Novelty:} we found no study that changes only the teacher across distinct base models.
\end{itemize}
Xuanyi Lyu runs training and sampling and leads the data-scaling experiment. Jiajun Chen builds the teacher, cross-style and human-control datasets and leads H3. Boshu Dou implements the metrics and the teacher classifier and leads the human evaluation. All of us annotate and write. The roughly 50 new runs need about 70 NVIDIA L4 GPU-hours, and the hosted teachers and rewriters need about 400K calls, which fits within our GCP credits. By the November 23 milestone we will have the 27 main-grid students with style normalization and first H1 and H2 results. By December 7 we will have the cross-style and data-scaling students and the human labels. The final report on December 17 will add teacher identification, robustness checks, qualitative analysis and any stretch goals.

\begingroup\small
\bibliographystyle{unsrt}
\bibliography{references}
\endgroup

\end{document}
